\scanpage{50}
$(s,x),(t,x)\in R$. On a

(i) $\forall s,t\in S\quad s\ne t$, existe au plus un $\unsure{a|x}\in E$
tel que $(s,x),(t,x)\in R$ i.e. $s,t\in R(x)$ i.e.
$\{s,t\}\subset R(x)$

\margin{\deleted{[\uline{NB} Cette condition est symétrique en $S,R$ elle signifie en fait \unsure{soit|si $s\ne t$} $x\ne y$ $(s,t\in S,\ x,y\in R)$ alors il existe au moins un $\{s,t\}$ \unsure{$s\mapsto x$}}}
\deleted{(ii) La relation dans $S$ :}

\uline{NB} Cette \deleted{relation}\inserted{condition} est symétrique en $S,R$. Elle
signifie que si $a\in\unsure{\mathcal P|P|\mathfrak P}_2(S)$, $u\in\mathcal P_2(R)$,
alors $\exists s\in a,\ \unsure{a|x}\in u$ tels que $\deleted{\unsure{R}}(s,x)\notin R$
% slanted
\margin{toute arête contenue dans \unsure{exactt} un $n$-uple}

Si on suppose que les \deleted{$R(x)$} $\operatorname{card}R(x)\geq2$
\deleted{\illegible{1}}\deleted{\unsure{alors}} (ce qui est le cas quand ça provient
d'une \deleted{\unsure{tri}} $n$-alité, \deleted{(cas \unsure{impropre})}\inserted{auquel cas $\operatorname{card}R(\unsure{n|x})=\deleted{\unsure{S}}n$}), alors
$R\longrightarrow\mathcal P(S)$ est injectif. De $=$, si on
suppose que pour tt $s\in S$, $ {}^{s}\!R(s)\subset R$ est
de cardinal $\geq2$ (ce qui est le cas si
ça provient d'une $n$-alité propre non triviale)
alors $S\longrightarrow\mathcal P(R)$ injectif)

(\unsure{ii|vi}) La relation dans $S$
\[
\{s,t\}\text{ n'est pas une arête}
\]
(i.e $s=t$ ou $\deleted{\unsure{\in|\nexists}}\nexists x\in E$ tel que $\{s,t\}\unsure{\subset|\not\subset}R(x)$)
est une relation d'équivalence dans $S$
i.e. pour \unsure{toute|une} arête $\{\unsure{r|s},t\}$, tout
sommet de $S$ est \unsure{liée|lié} soit à $r$, soit
à $t$.

Cette condition (ii) implique que $S=\coprod E_i$
de telle façon que \inserted{pour} $a\in\mathcal P_2(\unsure{R|S})$, $a$ est une
% archive sheet number: 49
